\documentclass[a4paper]{article}

% Basic packages
% Español (México) con XeLaTeX: fontspec para fuentes Unicode, polyglossia
% para la división silábica y los nombres del documento en la variante mexicana.
\usepackage{fontspec}
\usepackage{polyglossia}
\setdefaultlanguage[variant=mexican]{spanish}
% Los nombres chinos van como «pinyin (original)»: xeCJK da a los caracteres
% Han su propia fuente; sin ella XeLaTeX los omite con un aviso.
% FandolSong no cubre algunos caracteres de nombres (U+73FA); WenQuanYi Zen Hei sí.
\usepackage[AutoFallBack=true]{xeCJK}
\setCJKmainfont{FandolSong-Regular.otf}
\setCJKfallbackfamilyfont{\CJKrmdefault}{WenQuanYi Zen Hei}
% polyglossia no traduce los nombres de listings.
\AtBeginDocument{\providecommand{\lstlistingname}{}\renewcommand{\lstlistingname}{Listado}%
  \providecommand{\lstlistlistingname}{}\renewcommand{\lstlistlistingname}{Índice de listados}}
\usepackage{amsmath, amssymb}
\usepackage{graphicx}
\usepackage[margin=2.5cm]{geometry}
\usepackage[most]{tcolorbox}
\usepackage{etoolbox}
\usepackage{listings}
\usepackage{hyperref}
\usepackage{booktabs}
\usepackage{subcaption}
\usepackage{float}

% Highlight boxes
\newtcolorbox{knowledgebox}[1]{
    enhanced, colback=blue!5!white, colframe=blue!75!black, colbacktitle=blue!75!black,
    coltitle=white, fonttitle=\bfseries, title={#1}, attach boxed title to top left={yshift=-2mm, xshift=2mm},
    boxrule=1pt, sharp corners
}

\newtcolorbox{importantbox}[1]{
    enhanced, colback=yellow!10!white, colframe=yellow!80!black, colbacktitle=yellow!80!black,
    coltitle=black, fonttitle=\bfseries, title={#1}, sharp corners
}

\newtcolorbox{warningbox}[1]{
    enhanced, colback=red!5!white, colframe=red!75!black, colbacktitle=red!75!black,
    coltitle=white, fonttitle=\bfseries, title={#1}, sharp corners
}

% Listings
\lstset{
    language=Python,
    basicstyle=\ttfamily\small,
    keywordstyle=\color{blue},
    stringstyle=\color{red!60!black},
    commentstyle=\color{green!60!black},
    breaklines=true,
    frame=single,
    numbers=left,
    numberstyle=\tiny\color{gray},
    captionpos=b,
    extendedchars=true
}

% Front-page metadata
\newcommand{\notetitle}{Mesa redonda: el camino de los modelos grandes hacia la AGI}
\newcommand{\noteauthors}{Elaboradas a partir de materiales públicos del curso}
\newcommand{\notedate}{2024}
\newcommand{\videochannel}{Alibaba Cloud}
\newcommand{\videopublishdate}{2024}
\newcommand{\videoduration}{aproximadamente 60 minutos}
\newcommand{\videourl}{}
\newcommand{\videocoverpath}{cover.jpg}
\newcommand{\repourl}{https://github.com/hqhq1025/ai-course-notes}


% Footer with repo info
\usepackage{fancyhdr}
\pagestyle{fancy}
\fancyhf{}
\fancyfoot[C]{\thepage}
\fancyfoot[R]{\footnotesize\color{gray}\href{\repourl}{AI Course Notes}}
\renewcommand{\headrulewidth}{0pt}
\renewcommand{\footrulewidth}{0pt}

\begin{document}

\begin{titlepage}
\centering
{\Large Notas de la mesa redonda\par}
\vspace{1.2cm}
{\huge\bfseries \notetitle\par}
\vspace{0.8cm}
{\large \noteauthors\par}
\vspace{0.3cm}
{\large \notedate\par}
\vspace{1.2cm}

\ifdefempty{\videocoverpath}{
{\small Al generar las notas, indique la ruta local de la portada del video.\par}
}{
\includegraphics[width=0.82\textwidth,height=0.45\textheight,keepaspectratio]{\videocoverpath}\par
}

\vfill
\begin{tcolorbox}[width=0.9\textwidth, colback=black!2!white, colframe=black!60, sharp corners]
\textbf{Moderador}: Zhang Peng (张鹏) (fundador y presidente de GeekPark)\par
\textbf{Invitados}: Jiang Daxin (蒋大昕) (CEO de StepFun), Yang Zhilin (杨植麟) (fundador de Moonshot AI/Kimi), Zhu Jun (朱军) (científico jefe en Tsinghua University/Shengshu Technology)\par
\textbf{Canal}: \videochannel\par
\textbf{Fecha de publicación}: \videopublishdate\par
\textbf{Duración del video}: \videoduration\par
\textbf{Página del proyecto}: \href{\repourl}{\nolinkurl{\repourl}}\par
\end{tcolorbox}
\end{titlepage}

\tableofcontents
\newpage

%% --- Inicio del contenido --- %%

\section{Introducción: la aceleración de los últimos 18 meses}

Esta mesa redonda tuvo lugar poco después del lanzamiento de OpenAI O1. El moderador, Zhang Peng (张鹏), invitó a tres exploradores centrales del campo de los grandes modelos en China: Jiang Daxin (蒋大昕) de StepFun (阶跃星辰), Yang Zhilin (杨植麟) de Moonshot AI/Kimi (月之暗面/Kimi) y Zhu Jun (朱军) de Tsinghua (清华)/Shengshu Technology (生数科技), a repasar juntos el desarrollo de la AGI en los últimos 18 meses y a proyectar las direcciones futuras.

Los tres invitados coincidieron en que el desarrollo de la AGI se encuentra en un \textbf{estado de aceleración}, no de desaceleración.

\section{La doble aceleración de cantidad y calidad}

\subsection{De un solo dominante a la disputa entre varios}

Jiang Daxin (蒋大昕) resumió el progreso de los últimos 18 meses desde dos dimensiones: cantidad y calidad:

\begin{knowledgebox}{Dimensión de cantidad}
Cada mes surgen nuevos modelos y productos: Sora de OpenAI (febrero), GPT-4O (mayo), O1 (septiembre); la serie Claude 3 a 3.5 de Anthropic; la serie Gemini de Google; la serie Llama de Meta, entre otros. De un GPT-4 dominante en solitario se pasó a la disputa entre varios, persiguiéndose unos a otros.
\end{knowledgebox}

\begin{importantbox}{Dimensión de calidad: tres avances emblemáticos}
\begin{enumerate}
\item \textbf{GPT-4O}: un nuevo peldaño en la fusión multimodal, que unifica la comprensión visual (GPT-4V), la generación visual (DALL-E/Sora) y los modelos de voz (Whisper/Voice Engine). El mundo físico es en sí mismo multimodal, así que la fusión es una dirección inevitable.
\item \textbf{Tesla FSD V12}: un modelo grande de extremo a extremo que genera secuencias de control directamente a partir de señales de percepción. La conducción autónoma es un escenario emblemático del tránsito del mundo digital al mundo físico.
\item \textbf{O1}: la primera vez que se demuestra que un modelo de lenguaje también puede tener una capacidad de pensamiento lento (Sistema 2) como la del cerebro humano, lo cual es una condición básica para inducir el mundo.
\end{enumerate}
\end{importantbox}

\subsection{La doble expansión vertical y horizontal}

Yang Zhilin (杨植麟) analizó el asunto desde dos dimensiones:

\begin{itemize}
\item \textbf{Vertical (coeficiente intelectual)}: en matemáticas de competencia se pasó de reprobar por completo a más de 90 puntos; la capacidad de programación superó a programadores profesionales; la longitud de contexto saltó de 4--8K hasta 128K o incluso el nivel de un millón
\item \textbf{Horizontal (habilidades)}: la comprensión y generación multimodal, la conversión de artículos a podcast, la generación de video y otras transformaciones entre modalidades maduran cada vez más
\end{itemize}

\begin{knowledgebox}{La learning curve se está haciendo más pronunciada}
Zhu Jun (朱军) añadió una observación importante: la learning curve de cada campo se está haciendo cada vez más pronunciada. Por ejemplo, la generación de video pasó de que Sora sorprendiera al mundo a que la industria la alcanzara en solo medio año. La causa principal es que ya se cuenta con el conocimiento y la preparación sobre la ruta tecnológica, además de la madurez de la infraestructura en la nube.
\end{knowledgebox}

\subsection{Resumen de la sección}

En los últimos 18 meses, la AGI se aceleró notablemente tanto en cantidad (densidad de modelos) como en calidad (avances emblemáticos). El coeficiente intelectual vertical sigue subiendo, las habilidades horizontales se expanden con rapidez, y la velocidad con la que cada campo alcanza a los demás es cada vez mayor.
\section{O1: el nuevo paradigma de Scaling Law}

El lanzamiento de O1 fue el tema central de esta mesa redonda. Los tres invitados, desde ángulos distintos, dieron una evaluación positiva sumamente coincidente.

\subsection{La unificación de System 1 + System 2}

\begin{importantbox}{El análisis de Jiang Daxin (蒋大昕)}
El significado central de O1 es doble:
\begin{enumerate}
\item Demuestra por primera vez que un modelo de lenguaje puede tener capacidad de System 2: ya no es un pensamiento lineal (System 1), sino que puede explorar rutas distintas, reflexionar sobre sí mismo y autocorregirse, hasta encontrar la respuesta correcta
\item Combina el aprendizaje por imitación (preentrenamiento) con el aprendizaje por refuerzo, de modo que el modelo cuenta a la vez con capacidad de System 1 y de System 2
\end{enumerate}
\end{importantbox}

\begin{knowledgebox}{System 1 vs. System 2}
\textbf{System 1} (pensamiento rápido): la forma de pensar de GPT-4: aunque descomponga un problema complejo en varios pasos, sigue siendo una next-token prediction lineal.

\textbf{System 2} (pensamiento lento): la nueva capacidad de O1: puede explorar múltiples rutas, retroceder y corregir errores, hasta encontrar la solución correcta. Esta es la capacidad básica para inducir el mundo y resolver problemas nuevos.
\end{knowledgebox}

\subsection{RL Scaling: una nueva dimensión de Scaling}

Jiang Daxin (蒋大昕) señaló además que O1 trae un nuevo paradigma de \textbf{RL Scaling}:

\begin{itemize}
\item En el pasado, el aprendizaje por refuerzo de DeepMind (AlphaGo/AlphaFold/AlphaGeometry) siempre se diseñó para escenarios específicos
\item O1 logró por primera vez un RL Scaling a gran escala en cuanto a generalidad y capacidad de generalización
\item Todavía es solo un comienzo (preview), pero ya muestra una ruta con un límite superior muy alto
\end{itemize}

\subsection{El salto de nivel de la AGI}

\begin{knowledgebox}{La clasificación L1--L5 de AGI de Zhu Jun (朱军)}
\begin{itemize}
\item \textbf{L1 chatbot}: la era de ChatGPT
\item \textbf{L2 razonador}: la era de O1: en escenarios específicos (Narrow) ya se ha alcanzado el L2
\item \textbf{L3 agente}: pasar del mundo digital al mundo físico, para interactuar y transformar
\item \textbf{L4 innovador}: descubrir conocimiento nuevo, crear cosas nuevas
\item \textbf{L5 organizador}: coordinar organizaciones, operar con alta eficiencia
\end{itemize}
En cada nivel hay un salto de Narrow a General. O1 representa un cambio cualitativo notable de L1 a L2.
\end{knowledgebox}

\subsection{Resolver el Scaling Wall}

Yang Zhilin (杨植麟) señaló que O1 elevó el \textbf{límite superior} de la IA:

\begin{importantbox}{Del 5\% a las 10 veces}
Antes de O1, todos temían que el agotamiento de los datos (el muro de datos) hiciera fallar el Scaling Law. O1 demostró que, mediante el aprendizaje por refuerzo, se puede seguir haciendo Scaling: del límite superior de «aumentar la productividad en 5--10\%» a la posibilidad de «un PIB 10 veces mayor». En los setenta u ochenta años de historia de la IA, lo único que ha resultado eficaz es el Scaling, y O1 señaló una nueva dimensión para el Scaling.
\end{importantbox}

\subsection{Resumen de la sección}

El significado de O1: (1) el modelo de lenguaje obtiene capacidad de System 2; (2) el nuevo paradigma de RL Scaling rompe el muro de datos; (3) la AGI avanza de L1 a L2.

\section{El desafío de la generalización: de Narrow a General}

\subsection{El problema de la definición del Reward Model}

Zhu Jun (朱军) señaló el desafío central que enfrenta la generalización del RL:

\begin{warningbox}{El dilema de la definición del reward}
En la demostración de teoremas matemáticos y en la programación, el reward (función de recompensa) es inequívoco: una respuesta correcta es correcta, una incorrecta es incorrecta. Pero en dominios abiertos como la conducción autónoma, la creación artística o la generación de video, los límites entre ``bueno'' y ``malo'' son difusos, y la percepción de cada persona es distinta. Cómo definir el reward model en estos dominios es el problema técnico central de la generalización del RL.
\end{warningbox}

Además, obtener datos de supervisión de proceso (process supervision) también es difícil: se requiere etiquetar cada paso del proceso de razonamiento, lo cual exige personal especializado y datos de alto valor.

\subsection{La oportunidad para las empresas emergentes}

Yang Zhilin (杨植麟) analizó, desde la perspectiva de una empresa emergente, las oportunidades que trae el paradigma O1:

\begin{knowledgebox}{Oportunidades que trae la nueva variable tecnológica}
\begin{itemize}
\item Las empresas con cierto umbral de cómputo pueden hacer innovación fundamental a nivel de algoritmos de RL, e incluso lograr avances en el modelo base
\item Las empresas con menos cómputo pueden, mediante el post-entrenamiento, alcanzar lo mejor en un dominio específico
\item Tener una dirección definida pero un camino incierto es algo bueno para las empresas emergentes
\end{itemize}
\end{knowledgebox}

\subsection{Resumen de la sección}

El desafío central de la generalización del RL radica en la definición del reward en dominios abiertos y en la obtención de datos de supervisión de proceso. Pero el camino técnico ya está claro, y combinado con modelos base sólidos, la velocidad de generalización será mayor que la de la generación anterior de RL (la era de AlphaGo).

\section{Reacción en cadena en el panorama de la capacidad de cómputo}

\subsection{El reajuste del triángulo de la capacidad de cómputo}

Jiang Daxin (蒋大昕) analizó el efecto en cadena que el paradigma de RL tiene sobre el triángulo «algoritmo-capacidad de cómputo-datos»:

\begin{importantbox}{Tres juicios de certeza decreciente}
\begin{enumerate}
\item \textbf{Cierto}: la demanda de capacidad de cómputo del lado de la inferencia crece de manera multiplicativa (test-time scaling)
\item \textbf{Muy probablemente cierto}: la capacidad de cómputo de la etapa de entrenamiento con RL no es menor que la del preentrenamiento; los datos de Self-play no tienen, en teoría, un límite superior (por ejemplo, para entrenar Strawberry se usaron decenas de miles de tarjetas H100 durante varios meses)
\item \textbf{Incierto}: si la cantidad de parámetros del modelo principal necesita seguir haciendo scale; si el RL produce un efecto de «amplificador», el ROI del scaling de parámetros podría volver a ser positivo
\end{enumerate}
\end{importantbox}

\begin{warningbox}{Si se cumple el tercer punto}
Si el efecto de amplificación del RL vuelve a hacer eficaz el scaling de parámetros, entonces el crecimiento de la capacidad de cómputo regresaría a la dimensión cuadrática de $\text{cantidad de parámetros} \times \text{cantidad de datos}$: la demanda sobre toda la infraestructura de cómputo sería explosiva.
\end{warningbox}

\subsection{¿Quién sentirá primero la reevaluación de la capacidad de cómputo?}
Aunque los tres panelistas parten de ángulos distintos, sus juicios llegan a conclusiones cercanas: la presión sobre la capacidad de cómputo en el futuro no se manifestará solo en los clústeres de entrenamiento, sino que se transmitirá cuesta arriba a través de los chips, las plataformas en la nube, los servicios de inferencia y las empresas de aplicaciones. Es decir, lo que trae el RL Scaling no es un aumento de costo puntual, sino que toda la cadena de la industria tiene que recalcular el ROI.

\begin{table}[H]
\centering
\begin{tabular}{p{3.2cm}p{5.2cm}p{5cm}}
\toprule
Rol & Presión directa & Posible respuesta \\
\midrule
Empresas de modelos base & El entrenamiento y el test-time scaling elevan a la vez el presupuesto de GPU & Asignar con más precisión el presupuesto de preentrenamiento, posentrenamiento e inferencia \\
Proveedores de infraestructura en la nube & La demanda pico aumenta y la programación de los clústeres se vuelve más compleja & Ampliar el suministro de tarjetas de gama alta y mejorar el aprovechamiento multiinquilino \\
Empresas emergentes de aplicaciones & El costo de inferencia sube al llamar a modelos más potentes & Reducir el costo por solicitud mediante destilación, caché y flujos de trabajo especializados \\
\bottomrule
\end{tabular}
\caption{El RL Scaling afecta de manera distinta a los diferentes roles de la cadena de la industria}
\end{table}

\begin{knowledgebox}{Un juicio comercial implícito en la mesa redonda}
Si la demanda de capacidad de cómputo del lado de la inferencia estalla primero, quienes se beneficien primero no serán necesariamente los equipos con el modelo de más parámetros, sino posiblemente las empresas \textbf{más hábiles para «empaquetar la inferencia costosa como un servicio de alto valor»}. Por ello, el diseño de producto, la estrategia de caché y la división de flujos de trabajo serán tan importantes como la capacidad del modelo.
\end{knowledgebox}

\subsection{Resumen de la sección}
El cambio en el panorama de la capacidad de cómputo no es simplemente que «las tarjetas se encarecieron», sino que el RL Scaling vuelve a elevar tanto el entrenamiento como la inferencia, y transmite la presión hacia la nube, las aplicaciones y los modelos de negocio. Quien logre convertir una alta capacidad de cómputo en una entrega de alto valor tendrá más probabilidades de tomar la delantera en la siguiente ronda de competencia.

\section{Hoja de ruta de la evolución de la AGI}

\subsection{Simular el mundo $\rightarrow$ Explorar el mundo $\rightarrow$ Inducir el mundo}

Jiang Daxin (蒋大昕) propone una ruta de evolución de la AGI en tres etapas:

\begin{knowledgebox}{Marco de tres etapas}
\begin{enumerate}
\item \textbf{Simular el mundo}: representado por GPT-4O — fusión multimodal, para modelar el mundo físico
\item \textbf{Explorar el mundo}: representado por FSD V12 — control de extremo a extremo, del mundo digital al físico
\item \textbf{Inducir el mundo}: representado por O1 — pensamiento lento de Sistema 2, descubrimiento de patrones y conocimiento
\end{enumerate}
En los últimos 18 meses se lograron avances emblemáticos en las tres etapas.
\end{knowledgebox}

\subsection{Perspectiva histórica del Scaling}

Yang Zhilin (杨植麟) extrae una conclusión central de setenta u ochenta años de historia de la IA:

\begin{importantbox}{Lo único que ha funcionado en la historia de la IA es el Scaling}
Ya sea la cantidad de parámetros, la cantidad de datos o la cantidad de cómputo, más grande siempre es mejor. O1 no es solo un cambio cuantitativo, sino cualitativo: abrió una nueva dimensión de Scaling (RL Scaling) y disipó la inquietud de «qué hacer cuando se agoten los datos». Más importante aún, los datos generados por Self-play en teoría no tienen límite superior.
\end{importantbox}

\subsection{Resumen de la sección}

La evolución de la AGI puede entenderse a partir de las tres etapas «simular-explorar-inducir». El Scaling sigue siendo el motor central; O1 abrió la nueva dimensión de RL Scaling y rompió el límite del muro de datos.

\section{Síntesis y ampliación}

Esta mesa redonda se desarrolló justo después del lanzamiento de O1; los tres invitados, desde distintas dimensiones (la perspectiva emprendedora, la perspectiva académica y la perspectiva técnica), emitieron un juicio sobre el estado actual y el futuro de la AGI. El consenso central es el siguiente:

\begin{enumerate}
\item \textbf{Aceleración, no desaceleración}: en los últimos 18 meses el desarrollo de la AGI se ha acelerado, con una doble mejora en cantidad y calidad
\item \textbf{O1 es un cambio de paradigma}: abre una nueva dimensión de RL Scaling, dota al modelo de capacidad de System 2 y impulsa a la AGI de L1 a L2
\item \textbf{La generalización es el reto central}: pasar de Narrow a General requiere resolver los problemas de la definición de reward y de los datos de supervisión de proceso
\item \textbf{La demanda de cómputo estallará}: el lado de inferencia crecerá con certeza, el lado de entrenamiento crecerá con alta probabilidad y el número de parámetros podría reiniciar el Scaling
\item \textbf{Las oportunidades de emprendimiento se amplían}: hay una dirección definida (RL Scaling) pero una ruta incierta, lo que favorece a los equipos innovadores
\item \textbf{Ruta de tres etapas}: simular el mundo, explorar el mundo, inducir el mundo; en los últimos 18 meses ha habido avances emblemáticos en cada etapa
\end{enumerate}

\begin{knowledgebox}{Lecturas adicionales}
\begin{itemize}
\item Blog técnico de OpenAI O1: «Learning to Reason with LLMs», trabajo pionero del test-time scaling
\item Definición de niveles de AGI de OpenAI (L1--L5)
\item Análisis técnico de Tesla FSD V12: modelo grande de conducción autónoma de extremo a extremo
\item AlphaGo $\rightarrow$ AlphaFold $\rightarrow$ AlphaGeometry: la ruta de evolución de RL de DeepMind
\item Kahneman, «Thinking, Fast and Slow»: la obra original de la teoría de System 1/System 2
\end{itemize}
\end{knowledgebox}

%% --- Fin del contenido --- %%

\end{document}
